% =========================================================================
% تقرير مشروع مقرر معالجة صور عملي - المستوى الرابع
% جامعة إب - كلية الحاسبات والعلوم التطبيقية - قسم علوم الحاسوب وتقنية المعلومات
% العنوان: Pose2Skill-Robot: نظام ذكي لتعلّم ونقل المهارات البشرية إلى الروبوت باستخدام تقدير الوضعية والتفاعل مع الأجسام
% =========================================================================

\documentclass[11pt,a4paper]{article}

\usepackage{fontspec}
\usepackage{polyglossia}
\setmainlanguage{arabic}
\setotherlanguage{english}
\setmainfont{Amiri}

\usepackage{amsmath,amssymb}
\usepackage{geometry}
\usepackage{booktabs}
\usepackage{hyperref}
\usepackage{cite}

\geometry{a4paper, margin=2.2cm}

\begin{document}

% -------------------------------------------------------------
% صفحة الغلاف الأكاديمية الرسمية
% -------------------------------------------------------------
\begin{titlepage}
    \centering
    \vspace*{0.2cm}
    
    {\large \textbf{الجمهورية اليمنية}} \\
    \vspace{0.15cm}
    {\large \textbf{وزارة التعليم العالي والبحث العلمي}} \\
    \vspace{0.15cm}
    {\Large \textbf{جامعة إب}} \\
    \vspace{0.15cm}
    {\large \textbf{كلية الحاسبات والعلوم التطبيقية}} \\
    \vspace{0.15cm}
    {\normalsize \textbf{قسم علوم الحاسوب وتقنية المعلومات}} \\
    \vspace{0.15cm}
    {\small المستوى الرابع --- مقرر معالجة صور (عملي)} \\
    
    \vspace{2.2cm}
    \rule{\linewidth}{1.2pt} \\
    \vspace{0.5cm}
    {\large \textbf{تقرير مشروع مقرر معالجة صور عملي}} \\
    \vspace{0.4cm}
    {\huge \textbf{Pose2Skill-Robot: نظام ذكي لتعلّم ونقل المهارات البشرية إلى الروبوت باستخدام تقدير الوضعية والتفاعل مع الأجسام}} \\
    \vspace{0.4cm}
    {\large \textenglish{Pose2Skill-Robot: An Intelligent System for Learning and Transferring Human Skills to Robots Using Pose Estimation and Object Interaction}} \\
    \vspace{0.5cm}
    \rule{\linewidth}{1.2pt} \\
    
    \vspace{2.2cm}
    
    \begin{minipage}{0.48\textwidth}
        \begin{flushright} \large
            \textbf{إعداد الطلاب:} \\
            1. يعقوب المهاجري \\
            2. سليمان العربي \\
            3. مالك جبران
        \end{flushright}
    \end{minipage}
    \hfill
    \begin{minipage}{0.48\textwidth}
        \begin{flushleft} \large
            \textbf{إشراف أستاذ المقرر:} \\
            م. مالك المصنف \\
            \vspace{0.6cm}
            \small \textit{مشروع مقرر عملي}
        \end{flushleft}
    \end{minipage}
    
    \vfill
    
    {\large \textbf{العام الجامعي: 2025 -- 2026 م}} \\
    \vspace{0.4cm}
\end{titlepage}

\newpage

% -------------------------------------------------------------
% متن التقرير الأكاديمي
% -------------------------------------------------------------

\begin{abstract}
يقدم هذا التقرير التوثيق الفني والأكاديمي لمشروع مقرر معالجة صور عملي بقسم علوم الحاسوب وتقنية المعلومات بكلية الحاسبات والعلوم التطبيقية --- جامعة إب، بعنوان: \textbf{«Pose2Skill-Robot: نظام ذكي لتعلّم ونقل المهارات البشرية إلى الروبوت باستخدام تقدير الوضعية والتفاعل مع الأجسام»}. يهدف المشروع إلى تطبيق خوارزميات معالجة الصور الرقمية لتقدير وضعية الجسم البشري في الزمن الحقيقي وتوجيه روبوت بشري محاكى (SoftBank NAO) داخل بيئة Webots باستخدام كاميرا ويب أحادية دون الحاجة لأي حساسات ملبوسة. يعتمد النظام على استخراج 33 نقطة مفصلية للهيكل العظمي البشري عبر نموذج MediaPipe Pose، يتبعها اشتقاق الزوايا الكينماتيكية لمفاصل الأطراف العلوية والسفلية، مع تطبيق مرشح التنعيم الأسي (EMA) للحد من ارتعاش الصورة. كما يعالج المشروع إشكالية سقوط الروبوت الفيزيائية في بيئة المحاكاة عند رفع القدم للركل عبر حزام توازن برمجي مشرف (Supervisor Harness) يحافظ على استقرار الجذع مع ترك المفاصل حرة الحركة. ويتضمن النظام وحدة لتسجيل الحركات كمهارات مستقلة (JSON) وإعادة تشغيلها بدقة زمنية متطابقة عبر الاستيفاء الخطي (LERP).
\end{abstract}

\section{المقدمة وأهداف المشروع}
يعد استخراج الوضعيات الحركية للإنسان وتوجيه الروبوتات من التطبيقات البارزة لمقرر معالجة الصور والرؤية الحاسوبية. هدف هذا المشروع العملي إلى تطبيق خوارزميات معالجة صور الفيديو واستخلاص البيانات الحركية ونقلها لحظياً إلى روبوت بشري داخل بيئة محاكاة فيزيائية واقعية.

تم تحديد الأهداف الهندسية والعملية الواقعية التالية:
\begin{enumerate}
    \item قراءة ومعالجة تدفق الفيديو اللحظي من كاميرا الويب بكفاءة ومعدل إطارات مستقر (30 إطار/ثانية).
    \item استخراج إحداثيات مفاصل الجسم بدقة وتصفية الضجيج والاهتزازات الناتجة عن الكاميرا.
    \item حساب الزوايا المفصلية للأطراف العلوية (الأكتاف والأكواع) والسفلية (الأوراك والركبتين والكواحل) ومطابقتها مع مفاصل الروبوت NAO.
    \item معالجة مشكلة توازن الروبوت وسقوطه في بيئة المحاكاة عبر متحكم مشرف (Supervisor).
    \item توفير آلية لتسجيل الحركات المفيدة واستعادتها ذاتياً مع واجهة تفاعلية سهلة الاستخدام تعتمد على الأزرار الرسومية بالماوس.
\end{enumerate}

\section{الأساس الرياضي والمعالجة الكينماتيكية}

\subsection{حساب زوايا المفاصل ثلاثية الأبعاد}
لحساب زاوية المفصل $\theta$ المحصورة بين ثلاثة معالم مستخرجة من الكاميرا (مثل: الكتف $A$، الكوع $B$، المعصم $C$):
\begin{equation}
\vec{u} = \mathbf{p}_A - \mathbf{p}_B, \quad \vec{v} = \mathbf{p}_C - \mathbf{p}_B
\end{equation}
يتم استخراج زاوية الانحناء عبر الضرب القياسي للمتجهين:
\begin{equation}
\theta = \arccos\left( \frac{\vec{u} \cdot \vec{v}}{\|\vec{u}\| \|\vec{v}\| + \epsilon} \right)
\end{equation}
حيث تمثل $\epsilon = 10^{-7}$ حداً رياضياً لمنع القسمة على الصفر في حال تطابق النقاط.

\subsection{تنعيم الإشارات وتصفية الاهتزاز (EMA Filter)}
للحد من الارتعاش اللحظي الناتج عن تغيرات الإضاءة والتقدير البصري للكاميرا، تم تطبيق مرشح المتوسط المتحرك الأسي:
\begin{equation}
\hat{q}(t) = \alpha \cdot q(t) + (1 - \alpha) \cdot \hat{q}(t - \Delta t)
\end{equation}
تم اختيار المعامل $\alpha = 0.35$ أثناء التوجيه المباشر، مما وفر توازناً ملائماً بين الاستجابة السريعة واستقرار حركة مفاصل الروبوت. وعند تشغيل المهارات المسجلة يتم اعتماد $\alpha = 1.0$ لبلوغ المدى الحركي الأصلي بالكامل.

\subsection{مشغل الاستيفاء الخطي الزمني (Timestamp LERP)}
لتفادي ارتباط سرعة استعادة الحركة بعدد مرات تكرار الحلقة البرمجية أو تقلبات الكاميرا، تم اعتماد الاستيفاء الخطي المعتمد على الطوابع الزمنية الفعلية:
\begin{equation}
q(\tau) = (1 - \beta) q(t_k) + \beta q(t_{k+1}), \quad \beta = \frac{\tau - t_k}{t_{k+1} - t_k}
\end{equation}

\section{معمارية النظام والتنفيذ البرمجي}
يتكون النظام من المراحل المتسلسلة التالية:
\begin{itemize}
    \item \textbf{طبقة معالجة الصور:} استخدام واجهة DirectShow مع ترميز MJPG في OpenCV لتفادي تأخير الإطارات في نظام ويندوز.
    \item \textbf{نموذج التقدير البصري:} نموذج MediaPipe Pose لاستخراج 33 نقطة ثلاثية الأبعاد لكامل الجذع والأطراف.
    \item \textbf{حزام التوازن المشرف (Anti-Fall Harness):} برمجة متحكم الروبوت في Webots بصلاحيات المشرف، حيث يتم تثبيت السرعة المتجهية الخطية للجذع عبر البرمجة لمنع انهيار الروبوت عند رفع القدم، دون التأثير على سرعة محركات المفاصل.
    \item \textbf{محرك المهارات (LfD Engine):} حفظ المسارات الزمنية بصيغة JSON، وتشغيلها باستخدام الاستيفاء الخطي لضمان مطابقة سرعة الحركة الأصلية.
    \item \textbf{قناة الاتصال الشبكي:} إرسال البيانات اللحظية عبر مقابس TCP محلية (Port 10006) بهيكل بيانات خفيف الحجم.
\end{itemize}

\section{حل مشكلة التوازن الفيزيائي (Supervisor Harness)}
في المحاكاة الواقعية، يؤدي رفع قدم واحدة أو انحناء الروبوت السريع إلى خروج مركز الثقل عن قاعدة الارتكاز مما يسبب سقوط الروبوت فوراً. والحلول الشائعة القائمة على استدعاء إعادة تعيين الفيزياء كانت تؤدي إلى تصفير عزم المحركات وتجميد حركة الأذرع تماماً.

\textbf{الحل المعتمد في المشروع:}
\begin{enumerate}
    \item تشغيل متحكم الروبوت بصلاحية المشرف (\textenglish{Supervisor}).
    \item تثبيت سرعة حركة الجذع الخطية والدورانية حصرياً عبر البرمجة: \textenglish{torso.setVelocity([0, 0, 0, 0, 0, 0])}، مما يمتص قوى الانقلاب ويمنع السقوط دون المساس بسرعة أو دوران مفاصل الروبوت، فتتحرك الأرجل بحرية كاملة.
    \item برمجة آلية استعادة وقوف فورية تلقائية في حال حدوث أي ميلان يتجاوز 30 درجة.
\end{enumerate}

\section{مكتبة المهارات الحركية والواجهة الرسومية}
يتيح النظام للمستخدم تسجيل أي حركة جديدة كمهارة مستقلة بصيغة JSON عبر زر التسجيل المباشر على الشاشة. كما تم تزويد المشروع بمكتبة مهارات قياسية تم اختبارها والتحقق من دقتها:
\begin{itemize}
    \item \textbf{مهارة ركلة الفنون القتالية:} ركل بالقدم اليمنى مع موازنة الجسم.
    \item \textbf{مهارة القرفصاء الرياضي:} انحناء الركبتين والنزول لكامل الجسم والصعود بانتظام.
    \item \textbf{مهارة تحية النصر والذراعين:} رفع الذراعين للأعلى بتناسق.
    \item \textbf{مهارة التلويح والتحية:} رفع الذراع والتلويح للترحيب.
\end{itemize}

\section{النتائج التجريبية والقياسات}
أظهرت الاختبارات العملية كفاءة خط أنابيب المعالجة والاتصال وفق القياسات الموضحة في الجدول أدناه:

\begin{table}[htbp]
\centering
\caption{نتائج قياسات زمن المعالجة ومعدل الإطارات عبر مراحل النظام}
\label{tab:performance}
\begin{tabular}{lcc}
\toprule
\textbf{المرحلة البرمجية} & \textbf{متوسط زمن المعالجة} & \textbf{معدل الإطارات المحقق} \\
\midrule
التقاط الإطار من الكاميرا (MJPG 640x480) & 4.2 $\pm$ 0.8 مللي ثانية & 30 إطار/ثانية \\
استخراج معالم الجسم عبر MediaPipe & 16.5 $\pm$ 1.4 مللي ثانية & 30 إطار/ثانية \\
الحسابات الكينماتيكية وزوايا المفاصل & 0.6 $\pm$ 0.1 مللي ثانية & 30 إطار/ثانية \\
تصفية وتنعيم الإشارات (EMA) & 0.1 $\pm$ 0.05 مللي ثانية & 30 إطار/ثانية \\
الإرسال الشبكي عبر مقبس TCP المحلي & 1.1 $\pm$ 0.3 مللي ثانية & 30 إطار/ثانية \\
خطوة المحاكاة والفيزياء في Webots & 10.0 $\pm$ 0.5 مللي ثانية & 50 هرتز \\
\midrule
\textbf{التأخير الإجمالي (End-to-End Latency)} & \textbf{22.5 مللي ثانية} & \textbf{30 إطار/ثانية مستقر} \\
\bottomrule
\end{tabular}
\end{table}

\section{القيود الفنية والتطلعات المستقبلية}
بناءً على التجربة العملية والتقييم الهندسي الصادق للنظام، تبرز القيود الفنية والتحديات التالية:
\begin{enumerate}
    \item \textbf{حدود تتبع الأصابع:} يعتمد النظام على نموذج MediaPipe Pose العام المخصص لنقاط الهيكل العظمي للجسم (33 نقطة تنتهي عند المعصم ونقطة الكف الإجمالية)، ولا يتضمن تتبعاً مفصلياً لأصابع اليد الفردية، حيث يتطلب ذلك دمج نموذج تخصصي منفصل (مثل MediaPipe Hands ذو الـ 21 نقطة لكل يد).
    \item \textbf{التفاعل الفيزيائي وإمساك الأجسام:} إن مفاصل كف الروبوت NAO في المحاكاة هي سلاميات سلبية ذات احتكاك ضعيف، والتقاط الأشياء الصلبة وحملها حركياً يتطلب تطبيق خوارزميات متقدمة لتخطيط القبضة (Grasp Planning) وردود فعل لمسية (Tactile Feedback)، وهو ما يُمثل التطلع المستقبلي لتطوير بيئة التفاعل مع الأجسام.
\end{enumerate}

\section{الخاتمة}
حقق المشروع الأهداف التعليمية والتطبيقية لمقرر معالجة الصور العملي بنجاح، حيث تم ربط تقنيات تقدير الوضعية البصرية بالتحكم الكينماتيكي للروبوت وتثبيت توازنه الفيزيائي بدقة واستقرار وبمعدل إطارات حقيقي ومستدام.

\section*{المراجع (References)}
\begin{enumerate}
    \item Argall, B. D., Chernova, S., Veloso, M., \& Browning, B. (2009). A survey of robot learning from demonstration. \textit{Robotics and Autonomous Systems}, 57(5), 469–483.
    \item Lugaresi, C. et al. (2019). MediaPipe: A framework for building perception pipelines. \textit{arXiv preprint arXiv:1906.08172}.
    \item Michel, O. (2004). Cyberbotics Ltd. Webots: Professional mobile robot simulation. \textit{International Journal of Advanced Robotic Systems}, 1(1), 39–42.
    \item Schaal, S. (1999). Is imitation learning the route to humanoid robots? \textit{Trends in Cognitive Sciences}, 3(6), 233–242.
\end{enumerate}

\end{document}
